\section{Kinetic amplification for general homogeneous dispersion}
\label{sec:r13-general-dispersion}

The exterior kinetic term used for the biharmonic equation survives for
arbitrary positive elliptic homogeneous dispersion. The relevant density
is $|a(D)^{1/2}u|^2$. Anisotropic terms need not have a favorable sign:
after an explicit factorization they occur in a transition annulus, which
can be removed by averaging disjoint annuli.

Fix $m>2$, $s=m/2$, $N\ge1$, and a real even symbol $a$, homogeneous
of degree $m$, of class $C^{N+2}$ away from zero, satisfying
\[
 c_a|\xi|^m\le a(\xi)\le C_a|\xi|^m.
\]
Consider
\begin{equation}\label{r13g:equation}
 iu_t=a(D)u-|u|^{2m/N}u,
 \qquad p=2+2m/N,
 \qquad E(u)=\frac{A(u)}2-\frac{\|u\|_p^p}{p},
\end{equation}
where $A(u)=\|a(D)^{1/2}u\|_2^2$, $M(u)=\|u\|_2^2$, and
$P(u)=\operatorname{Im}\int\overline u\nabla u$.
The sharp critical mass $M_c$ is defined by
\[
 \|u\|_p^p\le\frac p2\left(\frac{M(u)}{M_c}\right)^{m/N}A(u).
\]
All dynamical statements refer to strong conservative $H^s$ solutions
for which the integrated localized identities pass by strong approximation.
Local well-posedness for every symbol and dimension is not asserted.

\begin{theorem}[Homogeneous kinetic amplification]
\label{thm:r13g-amplification}
There exist $d_0,C>0$, depending on $a,N,m$, such that every
forward-global solution satisfying
\[
 P=0,\quad E=-e<0,\quad
 d=\left(\frac{M}{M_c}\right)^{m/N}-1\in(0,d_0),\quad
 B=\sup_{t\ge0}A(u(t))<\infty
\]
obeys
\begin{equation}\label{r13g:amplification}
 e\le C d^{\alpha_m}B,
 \qquad \alpha_m=1+\frac{m}{2(m-1)}.
\end{equation}
Consequently, in a maximal strong flow with the energy-space
continuation alternative,
\begin{equation}\label{r13g:maximal}
 \sup_{0\le t<T_+} A(u(t))\ge C^{-1}|E|d^{-\alpha_m}.
\end{equation}
No radial symmetry, spatial moment, convexity of $a$, uniqueness of
the ground state, or real-sector linearized nondegeneracy is assumed.
\end{theorem}

\begin{corollary}[Finite kinetic excursions from standing waves]
\label{cor:r13g-ground-instability}
Let $Q$ be a real sharp standing-wave profile with $M(Q)=M_c$, and set
\[
 u_d(0)=(1+d)^{N/(2m)}Q.
\]
In the solution class of Theorem~\ref{thm:r13g-amplification},
\begin{equation}\label{r13g:ground-amplification}
 \sup_{0\le t<T_+(d)}
       \frac{A(u_d(t))}{A(u_d(0))}
 \ge c d^{-m/[2(m-1)]}.
\end{equation}
Every fixed finite amplification factor is therefore exceeded at a
finite forward time for all sufficiently small $d>0$. In particular
each such standing wave is orbitally unstable in $H^s$ modulo phase
and spatial translations.
\end{corollary}
\begin{proof}
Here $P=0$ and $E=-dA(u_d(0))/2$. Apply
\eqref{r13g:maximal}. A fixed amplification factor strictly below
the displayed lower bound is crossed at some time in the lifespan.
Since $u_d(0)\to Q$ in $H^s$, while $A^{1/2}$ is invariant under
phase and translation, a fixed sufficiently large amplification
separates the evolved state from the entire standing-wave orbit.
\end{proof}

For $m=4$ the exponent is $5/3$. The same proof covers positive
elliptic even homogeneous polynomials of degree $2k\ge4$, including
anisotropic ones, and $a(\xi)=|\xi|^m$ for every real $m>2$.
It does not exclude all bounded global negative-energy motion.

\subsection{An exterior kinetic term without convexity}

Put $b=\sqrt a$ and, for $\xi\ne0$, define the bounded degree-zero
symbols
\[
 B_{jk}(\xi)=\frac{\partial_j a(\xi)\xi_k}{a(\xi)}.
\]
Euler's identity gives $\sum_jB_{jj}=m$. These are real even symbols;
no sign is assigned to individual entries.

\begin{lemma}[Cutoff commutations]
\label{lem:r13g-cutoff}
For a fixed smooth compact cutoff $\theta$, $\theta_R(x)=\theta(x/R)$,
and $R\ge1$,
\begin{align}
 \|[B_{jk}(D),\theta_R]b(D)u\|_2
 &\le CR^{-1}\|u\|_{H^s},\label{r13g:zero-order-comm}\\
 \|[b(D),\theta_R]u\|_2
 &\le CR^{-1}\|u\|_{H^s}.\label{r13g:half-comm}
\end{align}
The same estimates hold for constant additions to $\theta_R$.
\end{lemma}
\begin{proof}
Write $h=\xi-\zeta$. The angular derivative bound for $B_{jk}$,
its boundedness, and $s>1$ give
\[
 |B_{jk}(\xi)-B_{jk}(\zeta)|b(\zeta)
 \le C|h||\zeta|^{s-1}.
\]
For $|h|\le|\zeta|/2$ use its homogeneous derivative of degree $-1$;
otherwise use $|\zeta|\le2|h|$. Young's convolution inequality and
$\||h|\widehat\theta_R(h)\|_1=O(R^{-1})$ prove
\eqref{r13g:zero-order-comm}. Similarly,
\[
 |b(\xi)-b(\zeta)|
 \le C|h|\bigl(|\zeta|+|h|\bigr)^{s-1}
\]
gives an upper bound
$CR^{-1}\|u\|_{H^{s-1}}+CR^{-s}\|u\|_2$ for
\eqref{r13g:half-comm}. Constants commute.
\end{proof}

Choose a radial $h_0\in C_c^\infty$ with $0\le h_0\le1$, equal
to one on $B_1$ and zero outside $B_2$. Let
\[
 f_R(x)=x h_0(x/R),\qquad
 L_R=f_R\cdot\nabla+\tfrac12\operatorname{div}f_R,
 \qquad C_R=[a(D),L_R].
\]

\begin{lemma}[Factorization of the principal form]
\label{lem:r13g-principal}
For $g=b(D)u$ and $R\ge1$,
\begin{equation}\label{r13g:factor}
 \langle u,C_Ru\rangle
 =\operatorname{Re}\sum_{j,k}
    \langle g,(\partial_j f_{R,k})B_{jk}(D)g\rangle
       +O(R^{-1}\|u\|_{H^s}^2).
\end{equation}
In particular,
\begin{align}
 \langle u,C_Ru\rangle
 &\le mA(u)-m\int_{|x|>2R}|g|^2
    +C\int_{R/2<|x|<4R}|g|^2
    +CR^{-1}\|u\|_{H^s}^2.
 \label{r13g:exterior-sign}
\end{align}
All constants are unchanged by translation of the cutoffs.
\end{lemma}
\begin{proof}
The Fourier kernel of $C_R$, with $h=\xi-\zeta$, is
\[
 \frac i2\,[a(\xi)-a(\zeta)](\xi+\zeta)_k
                                  \widehat f_{R,k}(h).
\]
The kernel of $\sum_{j,k}(\partial_j f_{R,k})q_{jk}(D)$, where
$q_{jk}(\zeta)=\partial_j a(\zeta)\zeta_k$, is
\[
 i h_j\widehat f_{R,k}(h)\partial_j a(\zeta)\zeta_k.
\]
Their difference is bounded by
\begin{equation}\label{r13g:kernel-remainder}
 C|h|^2(|\zeta|+|h|)^{m-1}|\widehat f_R(h)|.
\end{equation}
This is second-order Taylor expansion: the second derivatives of $a$
extend continuously at zero because $m>2$. To bound the form on $H^s$,
divide the kernel by $\langle\xi\rangle^s\langle\zeta\rangle^s$.
The elementary bound
\[
 \frac{(|\zeta|+|h|)^{m-1}}
      {\langle\xi\rangle^s\langle\zeta\rangle^s}
 \le C\langle h\rangle^{m-1}
\]
follows by splitting $|\zeta|\le2|h|$ and its complement.
Since $\widehat f_R(h)=R^{N+1}\widehat f_1(Rh)$, Schur's test
bounds \eqref{r13g:kernel-remainder} by $CR^{-1}$ as a form
$H^s\times H^s\to\mathbb C$.

Replacing $q_{jk}(D)$ by the $b(D)$ sandwich changes its kernel by
\[
 [b(\zeta)-b(\xi)]b(\zeta)B_{jk}(\zeta)
                      \widehat{\partial_j f_{R,k}}(h).
\]
Its magnitude is at most
$C|h|(|\zeta|+|h|)^{m-1}
|\widehat{\partial_j f_{R,k}}(h)|$.
The same weighted Schur estimate, now using the scale-invariant
coefficient $\partial_j f_{R,k}$, is $O(R^{-1})$.
Taking real parts proves \eqref{r13g:factor}.

Write
\[
 \partial_j f_{R,k}=h_0(x/R)\delta_{jk}+c_{jk,R}.
\]
The coefficients $c_{jk,R}$ are uniformly bounded and supported on
$R<|x|<2R$. The identity $\sum_jB_{jj}=m$ makes the diagonal part
exactly $m\int h_0(x/R)|g|^2$. Choose $\theta_R$ supported on
$R/2<|x|<4R$ and equal to one on the support of every $c_{jk,R}$.
For $c=c_{jk,R}$ and $B=B_{jk}(D)$,
\[
 \langle g,cBg\rangle
 =\langle\theta_Rg,cB(\theta_Rg)\rangle
   +\langle\theta_Rg,c[B,1-\theta_R]g\rangle.
\]
The first term is bounded by $C\|\theta_Rg\|_2^2$, and
Lemma~\ref{lem:r13g-cutoff} bounds the second by
$CR^{-1}\|u\|_{H^s}^2$.
Finally $h_0\le1$ and $h_0=0$ outside $2R$ give
\eqref{r13g:exterior-sign}. No pointwise principal-symbol sign has
been used.
\end{proof}

\subsection{Concentration and slow density transport}

The real two-component sharp polarization argument gives, for every
nonzero scalar complex field,
\begin{equation}\label{r13g:phase}
 \inf_\theta\left[
 \frac{M(\operatorname{Im}(e^{-i\theta}u))}{M(u)}+
 \frac{A(\operatorname{Im}(e^{-i\theta}u))}{A(u)}\right]
 \le C\left[1-
 \frac{\|u\|_p^p}{(p/2)(M(u)/M_c)^{m/N}A(u)}\right].
\end{equation}
For clarity, this applies the real $O(2)$ statement to
$(\operatorname{Re}u,\operatorname{Im}u)$. Its proof normalizes
$M=M_c,A=1$, uses compactness of sharp optimizers modulo translation
and phase, and writes $u=X+iY$ with $\langle Q,Y\rangle=0$.
Sharp interpolation for $X$ and transverse Taylor expansion show that
the deficit controls
\[
 \left\langle Y,
 \left[a(D)+\frac{m/N}{M_c}-|Q|^{2m/N}\right]Y\right\rangle
       -o(\|Y\|_{H^s}^2).
\]
The phase form has kernel $\mathbb RQ$ and is coercive on $Q^\perp$,
by the sharp density strict-convexity argument. Compactness of the
normalized optimizer family gives a uniform constant. Real-sector
nondegeneracy is not involved.

\begin{proof}[Proof of Theorem~\ref{thm:r13g-amplification}]
Critical scaling first reduces $B$ to one. Then
\[
 \frac{2e}{d}\le A(u(t))\le1,\qquad
 \|u(t)\|_{H^s}\le C.
\]
Equation~\eqref{r13g:phase} gives at each time a constant phase with
\[
 u=e^{i\theta}(X+iY),\qquad
 \|X\|_{H^s}\le C,\quad \|Y\|_{H^s}\le C\sqrt d.
\]
No time derivative of $\theta$ is used. For any real cutoff
$\chi_R(x)=\chi((x-z)/R)$, the Fourier estimate
\[
 |a(\xi)-a(\zeta)|
 \le C|\xi-\zeta|(|\zeta|+|\xi-\zeta|)^{m-1}
\]
and the weighted Schur argument above give
\[
 |\langle v,[a(D),\chi_R]w\rangle|
 \le CR^{-1}\|v\|_{H^s}\|w\|_{H^s}.
\]
The real-real and imaginary-imaginary terms make no contribution to
the mass derivative, since $a$ is real even. The nonlinear term
cancels pointwise. Thus
\begin{equation}\label{r13g:slow-mass}
 \left|\frac d{dt}\int\chi((x-z)/R)|u|^2\right|
 \le C\sqrt d/R.
\end{equation}

Sharp near-optimizer compactness and the lower bound $A\ge2e/d$
give a finite radius $R_0\ge1$ such that a ball of radius $R_0$
captures $M_c-\epsilon_0$ mass at each time. Here $\epsilon_0>0$
is a fixed sufficiently small fraction of $M_c$; $R_0$ may depend
on $e/d$. If two chosen centers at times separated by at most
$cR_0/\sqrt d$ were separated by more than $4R_0$, a cutoff mass
would change by a fixed positive amount, contradicting
\eqref{r13g:slow-mass}. Piecewise linear interpolation of centers
therefore supplies a Lipschitz center $a(t)$ with
\begin{equation}\label{r13g:center-speed}
 \|a'\|_\infty\le L=C\sqrt d,
\end{equation}
about which a fixed-radius ball contains $M_c-\epsilon_0$ mass
at all times. The speed constant is independent of $R_0$ and $e/d$.

For every $\eta>0$, some finite $R_\eta$ satisfies
\begin{equation}\label{r13g:tail-mass}
 \sup_t\int_{|x-a(t)|>R_\eta}|u(t,x)|^2
 \le M-M_c+\eta=:\mu.
\end{equation}
Indeed a violating sequence, translated by $a(t)$, has a weak $H^s$
limit $v$ with $M_c-\epsilon_0\le M(v)\le M_c-\eta$.
Its remainder also has mass strictly below $M_c$, for small $d_0$.
Quadratic kinetic splitting and Brezis--Lieb splitting would then
express $-e$ as a limit of two nonnegative subcritical energies.
This contradiction proves \eqref{r13g:tail-mass}.

\subsubsection*{Moving virial and exterior absorption}
Let $g=a(D)^{1/2}u$ and define
\[
 J_R=2\operatorname{Im}\int f_R(x-a(t))\cdot\overline u\nabla u,
\quad
 K_R=\int_{|x-a|>8R}|g|^2,
\quad
 H_R=\int_{R/4<|x-a|<8R}|g|^2.
\]
Since $s>1$, $|J_R|\le CR$. The integrated localized identity,
Lemma~\ref{lem:r13g-principal}, and $P=0$ give, for $R\ge4R_\eta$,
\begin{align}
 J_R'\le{}&-4me-2mK_R+CH_R+CR^{-1}
  +C\int_{|x-a|>R}|u|^p\notag\\
 &+CL\int_{|x-a|>R}|u||\nabla u|.
 \label{r13g:virial}
\end{align}
The local nonlinear term is
$-2(p-2)p^{-1}\int\operatorname{div}f_R|u|^p$; its difference
from the full dilation expression is supported outside $R$.
Motion of the center contributes $-2a'\cdot P$ plus the displayed
exterior current, because $Df_R=I$ on $B_R$.

Choose $\zeta_R=0$ on $B_{R/4}$ and $\zeta_R=1$ outside
$B_{R/2}$. Equations~\eqref{r13g:half-comm} and
\eqref{r13g:tail-mass} imply
\[
 M(\zeta_Ru)\le\mu,\qquad
 A(\zeta_Ru)\le C(K_R+H_R)+CR^{-2}.
\]
Sharp interpolation and interpolation between derivative orders zero
and $s$ therefore give
\begin{align}
 \int_{|x-a|>R}|u|^p
 &\le C\mu^{m/N}(K_R+H_R+R^{-2}),\notag\\
 \int_{|x-a|>R}|u||\nabla u|
 &\le C\mu^{1-1/m}(K_R+H_R+R^{-2})^{1/m}.
 \label{r13g:tail-bounds}
\end{align}
Take $d_0$ and then $\eta$ sufficiently small. The $K_R$ term in
the first inequality is absorbed by its negative coefficient in
\eqref{r13g:virial}. Young's inequality absorbs the current's
$K_R^{1/m}$ term, leaving
\begin{equation}\label{r13g:averagable}
 J_R'\le -4me+CH_R
    +CL\mu^{1-1/m}H_R^{1/m}
    +CL^{m/(m-1)}\mu+o_R(1).
\end{equation}
The remainder tends to zero uniformly in time on the normalized
$H^s$ ball.

\subsubsection*{Annulus averaging and restoration of scale}
Fix $K<\infty$ and $R_*\ge4R_\eta$, and take
$R_j=64^jR_*$, $0\le j<K$. The annuli defining $H_{R_j}$ are
disjoint, so $\sum_jH_{R_j}(t)\le1$. For each time window $[0,T]$
one annulus has time average at most $1/K$. Integrating
\eqref{r13g:averagable}, using Jensen's inequality, and estimating
the endpoints of $J_R$ yields
\[
 4me\le \frac{C64^{K-1}R_*}{T}+\frac CK
  +CL\mu^{1-1/m}K^{-1/m}+CL^{m/(m-1)}\mu
  +\sup_{R\ge R_*}|o_R(1)|.
\]
Let $T\to\infty$, then $K\to\infty$, then $R_*\to\infty$,
and finally $\eta\downarrow0$. No time-dependent cutoff radius
has been differentiated. We obtain
\[
 e\le CL^{m/(m-1)}(M-M_c)
 \le C d^{1+m/[2(m-1)]}.
\]
Under $u_\lambda(t,x)=\lambda^{N/2}u(\lambda^m t,\lambda x)$,
$A$ and $E$ scale by $\lambda^m$ and $d$ stays fixed. Choosing
$\lambda=B^{-1/m}$ restores \eqref{r13g:amplification}.
If the maximal solution is finite-time or global unbounded, its
kinetic supremum is infinite by the continuation alternative; the
only remaining case is the bounded global one already proved.
\end{proof}

\begin{remark}[The remaining dynamical obstruction]
After annulus averaging the only surviving exterior cost is
$L^{m/(m-1)}(M-M_c)$. Zero momentum does not set $L$ to zero:
density transport involves $\nabla a(D)$, whereas momentum involves
$D$. The present proof gives $L=O(\sqrt d)$, sufficient for the
amplification law but insufficient to exclude every large bounded
global orbit for fixed $d>0$.
\end{remark}

\subsection{Homogeneous Hartree attraction}
\label{subsec:r13g-hartree}

The same exterior kinetic mechanism applies to a nonlocal interaction.
Assume now $2<m<N$, and let $K$ be real, even, $C^1$ away from zero,
homogeneous of degree $-m$, and positive definite as a density quadratic
form. Suppose
\[
 |K(x)|+|x||\nabla K(x)|\le C_K|x|^{-m}.
\]
Consider
\[
 iu_t=a(D)u-(K*|u|^2)u,\qquad
 H_K(u)=\iint K(x-y)|u(x)|^2|u(y)|^2\,dx\,dy,
 \qquad E=\frac A2-\frac{H_K}4.
\]
Let $M_c$ be the sharp mass in $H_K(u)\le2(M/M_c)A(u)$.
The positive density form supplies sharp compactness and phase
coercivity as in the scalar Hartree foundations.

\begin{proposition}[Hartree amplification]
\label{prop:r13g-hartree}
Theorem~\ref{thm:r13g-amplification} and its maximal-flow conclusion
hold for this Hartree equation with
\[
 d=M/M_c-1,\qquad \alpha_m=1+\frac{m}{2(m-1)}.
\]
For a real sharp ground state, the amplitude data
$u_d(0)=\sqrt{1+d}\,Q$ satisfy
\eqref{r13g:ground-amplification} in the same conservative strong
solution class.
\end{proposition}
\begin{proof}
The real two-component phase estimate uses the deficit
\[
 1-\frac{H_K(u)}{2(M/M_c)A(u)}
 =\frac{d-2e/A(u)}{1+d}\le d.
\]
Density transport is unchanged because the nonlinear potential is
real. Sharp compactness and Hartree energy splitting yield the same
center and surplus-mass bound. The splitting follows by local
compactness and decay of the potential of a smooth compact core;
Hardy--Littlewood--Sobolev controls its approximation error.

Only the nonlinear virial defect needs a new estimate. With
$\rho=|u|^2$, its symmetrized expression is
\[
 \iint\rho(x)\rho(y)
 [f_R(x-a)-f_R(y-a)]\cdot\nabla K(x-y)\,dx\,dy.
\]
Subtracting the full dilation value $-mH_K(u)$ gives an integrand
which vanishes when both points are in $B(a,R)$. Since $f_R$ is
uniformly Lipschitz, the absolute remaining integrand is bounded by
$C\rho(x)\rho(y)|x-y|^{-m}$. By symmetry suppose $|x-a|>R$.
If $|y-a|<R/2$, then $|x-y|\ge R/2$, and the contribution is bounded
by $CR^{-m}M^2$. Otherwise both points lie outside $B(a,R/2)$.
Choose $\zeta_R=1$ there and zero on $B(a,R/4)$. The remaining
contribution is at most the absolute-Riesz Hartree form of
$w=\zeta_Ru$. Hardy--Littlewood--Sobolev and interpolation give
\[
 \iint|x-y|^{-m}|w(x)|^2|w(y)|^2\,dx\,dy
 \le C\|w\|_{4N/(2N-m)}^4
 \le CM(w)A(w).
\]
Consequently the nonlinear defect satisfies
\begin{equation}\label{r13g:hartree-tail}
 |\text{Hartree defect}|
 \le CR^{-m}M^2+C\mu(K_R+H_R+R^{-2}).
\end{equation}
The term $C\mu K_R$ is absorbed for small $d_0,\eta$.
All other steps of the moving virial, Young absorption, and annulus
averaging are identical. Finally, amplitude data satisfy
$E=-dA(u_d(0))/2$, giving the stated excursion law.
\end{proof}

The spatial separation in \eqref{r13g:hartree-tail} is essential:
the true core--tail interaction decays with $R$, and the remaining
near-diagonal interaction sees only the small exterior mass. Estimating
one exterior factor against the entire state by HLS would instead
leave a square-root exterior-energy term and lose this absorption.
